\section{Mengen}
\label{sec:fundamentals_sets}

Der Mengenbegriff ermöglicht es uns später unter anderem,
Zahlenräume abzugrenzen und Definitions- sowie Wertemengen von Funktionen präzise anzugeben.

\begin{DefinitionBox}{Menge}{fundamentals_set}
  Eine \emph{Menge} \EnglishTerm{set} ist eine Zusammenfassung bestimmter
  \enquote{Elemente} \EnglishTerm{elements} zu einem Ganzen. Jedes Element
  kann höchstens einmal in einer Menge enthalten sein; die Reihenfolge der
  Elemente spielt keine Rolle.
\end{DefinitionBox}

Elemente einer Menge können beispielsweise Zahlen, Objekte oder selbst wieder Mengen sein.\\

Um eine Menge eindeutig anzugeben, werden zwei Schreibweisen verwendet:
\begin{enumerate}[label=(\roman*),itemsep=0.8em]
  \item \textbf{Aufzählende Form:} Eine Menge, die die Elemente $a$, $b$ und
    $c$ enthält, schreibt man als
    \begin{ClozeArea}[33]{2}
      \[
        M=\{a,b,c\}.
      \]
    \end{ClozeArea}
  \item \textbf{Beschreibende Form:} Eine Menge, die alle Elemente mit einer
    bestimmten Eigenschaft enthält, schreibt man als
    \begin{ClozeArea}[33]{2}
      \[
        M=\{x\mid x\text{ besitzt die Eigenschaft }E\}.
      \]
    \end{ClozeArea}
\end{enumerate}

\begin{ExampleBox}{Aufzählende vs. Beschreibende Form}{fundamentals_set_notation_example}
  Die Menge $M$ der ganzen Zahlen von $1$ bis $5$ kann auf zwei Arten
  geschrieben werden:
  \begin{align*}
    \text{Aufzählende Form:}\quad
      M &= \Cloze{6.3cm}{\{1,2,3,4,5\}},\\[0.5em]
    \text{Beschreibende Form:}\quad
      M &= \Cloze{6.3cm}{\{n\in\mathbb{N}\mid 1\leq n\leq 5\}}.
  \end{align*}
\end{ExampleBox}

Auch eine Menge ohne Elemente ist eine wohldefinierte Menge. Sie wird als
leere Menge bezeichnet.

\begin{DefinitionBox}{Leere Menge}{fundamentals_emptyset}
  Die \emph{leere Menge} \EnglishTerm{empty set} ist die Menge, die kein
  Element enthält. Sie kann gleichwertig mit dem Symbol $\emptyset$ oder durch
  leere geschweifte Klammern $\{\}$ (aufzählende Form) bezeichnet werden:
  \[
    \emptyset=\{\}
  \]
\end{DefinitionBox}



\begin{minipage}{\linewidth}
Um sinnvoll mit Mengen arbeiten zu können, müssen wir angeben können, ob ein
bestimmtes Objekt Element einer Menge ist oder nicht.

\begin{DefinitionBox}{Element einer Menge}{fundamentals_set_element}
  Ist $a$ ein \emph{Element} \EnglishTerm{element} aus der Menge $M$, so
  schreibt man
  \begin{align*}
    &\Cloze{2.8cm}{a\in M}
    \qquad \text{\enquote{ist Element von}}, \\
    %
    \text{sonst} \quad
    &\Cloze{2.8cm}{a\notin M}
    \qquad\text{\enquote{ist kein Element von}}.
  \end{align*}
\end{DefinitionBox}

\begin{ExampleBox}{Elemente von Mengen}{fundamentals_elements_of_sets}
  \begin{enumerate}[label=\alph*),itemsep=0.6em]
    \item $1\ \Cloze{1.5cm}{\in}\ \{1,2,3\}$
    \item $\text{Peter}\ \Cloze{1.5cm}{\notin}\
      \{\text{Franz},\text{Anne}\}$
  \end{enumerate}
\end{ExampleBox}
\end{minipage}



\subsection{Teilmengen}
\label{subsec:fundamentals_subsets}

\begin{DefinitionBox}{Teilmenge}{fundamentals_subset}
  \begin{minipage}[t]{0.62\linewidth}
    \vspace{0pt}
    Eine Menge $A$ heißt \emph{Teilmenge} \EnglishTerm{subset} einer Menge
    $B$, wenn jedes Element von $A$ auch Element von $B$ ist. Man schreibt
    \[
      A\subseteq B
      \quad\Longleftrightarrow\quad
      \forall x:\; x\in A\Rightarrow x\in B.
    \]
    Ist $A$ keine Teilmenge von $B$, existiert mindestens ein Element von $A$,
    das nicht in $B$ liegt:
    \[
      A\not\subseteq B
      \quad\Longleftrightarrow\quad
      \exists x:\; x\in A\land x\notin B.
    \]
  \end{minipage}\hfill
  \begin{minipage}[t]{0.32\linewidth}
    \vspace{0pt}
    \centering
    \resizebox{\linewidth}{!}{%

      \CachedTikzFigure{content/resources/chapters/fundamentals/tikz/subset_inclusion.tex}%
    }
  \end{minipage}
  \par\smallskip
  \hfill
  {\footnotesize\color{black!55}
    $\forall\mathrel{\widehat{=}}\text{„für alle“},\qquad
     \exists\mathrel{\widehat{=}}\text{„es existiert“}$.}
\end{DefinitionBox}

\begin{ExampleBox}[after skip=0pt]{Teilmengen}{fundamentals_subset_example}
  \begin{enumerate}[label=\alph*),itemsep=0.6em]
    \item $\{1,2,3\}\ \Cloze{1.8cm}{\subseteq}\ \{1,4,3,2\}$
    \item $\{1,2\}\ \Cloze{1.8cm}{\not\subseteq}\ \{1,3\}$,
      denn $\Cloze{2.8cm}{2\notin\{1,3\}}$.
    \item $\emptyset\ \Cloze{1.8cm}{\subseteq}\ M\quad(\forall M)$
    \par\smallskip
    Die leere Menge ist Teilmenge jeder Menge.
  \end{enumerate}
\end{ExampleBox}

\par\nobreak
\begin{RemarkBox}[remark title=Echte Teilmengen,before skip=0.35\baselineskip]
  Bei $A\subseteq B$ ist die Gleichheit $A=B$ erlaubt. Gilt zusätzlich
  $A\neq B$, so heißt $A$ eine \textbf{echte Teilmenge}
  \EnglishTerm{proper subset} von $B$; man schreibt $A\subsetneq B$.
\end{RemarkBox}

\clearpage
\subsection{Mengenoperationen}
\label{subsec:fundamentals_set_operations}

Zwischen Mengen existieren Operationen, mit denen aus gegebenen Mengen neue
Mengen gebildet werden können. Zu den wichtigsten gehören Schnitt, Vereinigung
und Differenz.

\subsubsection{Schnittmenge}
\label{subsubsec:fundamentals_set_intersection}

\noindent\begin{minipage}{\linewidth}
\begin{DefinitionBox}{Schnittmenge}{fundamentals_set_intersection}
  \begin{minipage}[c]{0.70\linewidth}
    Die \emph{Schnittmenge} \EnglishTerm{intersection} zweier Mengen $A$ und
    $B$ enthält alle Elemente, die sowohl in der Menge $A$ als auch in der
    Menge $B$ enthalten sind.
    \begin{ClozeArea}[24]{3}
      \[
        A\cap B=\{x\mid x\in A\land x\in B\}.
      \]
    \end{ClozeArea}
  \end{minipage}\hfill
  \begin{minipage}[c]{0.25\linewidth}
    \centering
    \CachedTikzFigure{content/resources/chapters/fundamentals/tikz/set_intersection.tex}
  \end{minipage}
\end{DefinitionBox}

\begin{ExampleBox}{Schnittmenge}{fundamentals_set_intersection_example}
  \begin{enumerate}[label=\alph*),itemsep=0.6em]
    \item $\{1,2\}\cap\{2,3\}=\Cloze{3cm}{\{2\}}$
    \item $\{1,2,3\}\cap\{1,2,3\}=\Cloze{3cm}{\{1,2,3\}}$
    \item $M\cap\emptyset=\Cloze{2.5cm}{\emptyset}\quad(\forall M)$
  \end{enumerate}
\end{ExampleBox}
\end{minipage}



\subsubsection{Vereinigungsmenge}
\label{subsubsec:fundamentals_set_union}

\noindent\begin{minipage}{\linewidth}
\begin{DefinitionBox}{Vereinigungsmenge}{fundamentals_set_union}
  \begin{minipage}[c]{0.70\linewidth}
    Die \emph{Vereinigungsmenge} \EnglishTerm{union} zweier Mengen $A$ und
    $B$ enthält alle Elemente, die in der Menge $A$, in der Menge $B$ oder in
    beiden Mengen enthalten sind.
    \begin{ClozeArea}[24]{3}
      \[
        A\cup B=\{x\mid x\in A\lor x\in B\}.
      \]
    \end{ClozeArea}
  \end{minipage}\hfill
  \begin{minipage}[c]{0.25\linewidth}
    \centering
    \CachedTikzFigure{content/resources/chapters/fundamentals/tikz/set_union.tex}
  \end{minipage}
\end{DefinitionBox}

\begin{ExampleBox}{Vereinigungsmenge}{fundamentals_set_union_example}
  \begin{enumerate}[label=\alph*),itemsep=0.6em]
    \item $\{1,2\}\cup\{2,3\}=\Cloze{3cm}{\{1,2,3\}}$
    \item $\{1,2,3\}\cup\{1,2,3\}=\Cloze{3cm}{\{1,2,3\}}$
    \item $M\cup\emptyset=\Cloze{2.5cm}{M}\quad(\forall M)$
  \end{enumerate}
\end{ExampleBox}
\end{minipage}


\subsubsection{Differenzmenge}
\label{subsubsec:fundamentals_set_difference}

\noindent\begin{minipage}{\linewidth}
\begin{DefinitionBox}{Differenzmenge}{fundamentals_set_difference}
  \begin{minipage}[c]{0.70\linewidth}
    Die \emph{Differenzmenge} \EnglishTerm{set difference} $A\setminus B$
    enthält alle Elemente, die in der Menge $A$, aber nicht in der Menge $B$
    enthalten sind.
    \begin{ClozeArea}[24]{3}
      \[
        A\setminus B=\{x\mid x\in A\land x\notin B\}.
      \]
    \end{ClozeArea}
  \end{minipage}\hfill
  \begin{minipage}[c]{0.25\linewidth}
    \centering
    \CachedTikzFigure{content/resources/chapters/fundamentals/tikz/set_difference.tex}
  \end{minipage}
\end{DefinitionBox}

\begin{ExampleBox}{Differenzmenge}{fundamentals_set_difference_example}
  \begin{enumerate}[label=\alph*),itemsep=0.6em]
    \item $\{1,2\}\setminus\{2,3\}=\Cloze{3cm}{\{1\}}$
    \item $\{1,2,3\}\setminus\{1,2,3\}=\Cloze{2.5cm}{\emptyset}$
    \item $M \setminus \emptyset = \Cloze{2.5cm}{M}$
  \end{enumerate}
\end{ExampleBox}
\end{minipage}

\clearpage
\subsection{Zahlenmengen}
\label{subsec:fundamentals_number_sets}

Ausgehend von den natürlichen Zahlen werden die Zahlenmengen schrittweise
erweitert. Der Grund ist stets derselbe: Eine Rechenoperation soll ausführbar
sein, ihr Ergebnis liegt aber nicht mehr in der bisherigen Zahlenmenge.

\begin{DefinitionBox}{Abgeschlossenheit}{fundamentals_closure}
  Eine Zahlenmenge $M$ heißt bezüglich einer Rechenoperation $\circ$
  \emph{abgeschlossen} \EnglishTerm{closed}, wenn für alle $a,b\in M$ auch
  \[
    a\circ b\in M
  \]
  gilt, sofern die Rechenoperation definiert ist.
\end{DefinitionBox}

\subsubsection{Natürliche Zahlen}
\label{subsubsec:fundamentals_natural_numbers}

Die \emph{natürlichen Zahlen} dienen dem Zählen:
\begin{align*}
  \N
    &= \Cloze{7cm}{\{1,2,3,4,\ldots\}},\\[0.4em]
  \N_0
    &= \Cloze{7cm}{\N\cup\{0\}=\{0,1,2,3,\ldots\}}.
\end{align*}
Die Addition und Multiplikation natürlicher Zahlen führen wieder zu einer
natürlichen Zahl. Die natürlichen Zahlen sind bezüglich $+$ und $\cdot$
abgeschlossen:
\[
  3+5=\Cloze{1.0cm}{8}\in\N
  \qquad
  3\cdot5=\Cloze{1.2cm}{15}\in\N
  \qquad \text{mit } 3\in\N,\ 5\in\N.
\]

Bezüglich der Subtraktion sind die natürlichen Zahlen dagegen nicht
abgeschlossen:
\[
  3-5=\Cloze{1.2cm}{-2}\notin\N
  \qquad \text{mit } 3\in\N,\ 5\in\N.
\]


\subsubsection{Ganze Zahlen}
\label{subsubsec:fundamentals_integer_numbers}

Damit auch solche Differenzen erfasst werden, ergänzt man die negativen Zahlen
und erhält die \emph{ganzen Zahlen}:
\[
  \Z=\Cloze{9cm}{\{\ldots,-3,-2,-1,0,1,2,3,\ldots\}}.
\]
\begin{center}
  \begingroup
    \def\NumberSetStage{2}
    \CachedTikzFigure{content/resources/chapters/fundamentals/tikz/number_set_inclusions.tex}
  \endgroup
\end{center}

In $\Z$ sind Addition, Subtraktion und Multiplikation abgeschlossen.

\[
  -3+5=\Cloze{1.0cm}{2}\in\Z
  \qquad
  -3-5=\Cloze{1.2cm}{-8}\in\Z
  \qquad
  (-3)\cdot5=\Cloze{1.3cm}{-15}\in\Z
  \qquad \text{mit } -3\in\Z,\ 5\in\Z.
\]

Bezüglich der Division sind die ganzen Zahlen nicht abgeschlossen:
\[
  -3:5=-\frac35=\Cloze{1.8cm}{-0{,}6}\notin\Z
  \qquad \text{mit } -3\in\Z,\ 5\in\Z.
\]


\subsubsection{Rationale Zahlen}
\label{subsubsec:fundamentals_rational_numbers}

Damit auch Quotienten ganzer Zahlen erfasst werden, erweitert man $\Z$ zu den
\emph{rationalen Zahlen}:
\[
  \Q=\left\{\frac pq\,\middle|\,p\in\Z,\ q\in\N\right\}
  \qquad \text{„Menge aller Brüche aus ganzen Zahlen“}.
\]
\begin{center}
  \begingroup
    \def\NumberSetStage{3}
    \CachedTikzFigure{content/resources/chapters/fundamentals/tikz/number_set_inclusions.tex}
  \endgroup
\end{center}
Jede rationale Zahl lässt sich als Bruch ganzer Zahlen darstellen. Endliche
und periodische Dezimalzahlen sind daher rational, zum Beispiel
\[
  -\frac12=\Cloze{1.8cm}{-0{,}5}\in\Q,
  \qquad
  \frac19=\Cloze{2.0cm}{0{,}\overline{1}}\in\Q.
\]
Die rationalen Zahlen sind bezüglich $+$, $-$, $\cdot$ und der Division durch
eine von null verschiedene Zahl abgeschlossen.

\[
  \begin{aligned}
    \frac12+\left(-\frac14\right)
      &=\Cloze{3.0cm}{\frac14}\in\Q,
    &\qquad
    \frac12-\left(-\frac14\right)
      &=\Cloze{3.0cm}{\frac34}\in\Q,\\[0.5em]
    \frac12\cdot\left(-\frac14\right)
      &=\Cloze{3.0cm}{-\frac18}\in\Q,
    &\qquad
    \frac12:\left(-\frac14\right)
      &=\Cloze{3.0cm}{-2}\in\Q
  \end{aligned}
  \qquad
  \text{mit } \frac12\in\Q,\ -\frac14\in\Q.
\]

Bezüglich des Wurzelziehens sind die rationalen Zahlen allerdings nicht
abgeschlossen. So gibt es keine rationale Zahl, deren Quadrat gleich $2$ ist:
\[
  x^2=2
  \quad\Longrightarrow\quad
  x=\pm\Cloze{1.2cm}{\sqrt2}\notin\Q.
\]


\subsubsection{Reelle Zahlen}
\label{subsubsec:fundamentals_real_numbers}

Zahlen, die sich nicht als Quotient ganzer Zahlen darstellen lassen, heißen
\emph{irrationale Zahlen}. Beispiele sind
\[
  \sqrt2=1{,}4142\ldots,
  \qquad
  \pi=3{,}14159\ldots,
  \qquad
  \mathrm e=2{,}71828\ldots.
\]
Rationale und irrationale Zahlen bilden zusammen die \emph{reellen Zahlen}.

\begin{center}
  \begingroup
    \def\NumberSetStage{4}
    \CachedTikzFigure{content/resources/chapters/fundamentals/tikz/number_set_inclusions.tex}
  \endgroup
\end{center}

Die reellen Zahlen $\R$ sind bezüglich Addition, Subtraktion und Multiplikation
sowie bezüglich der Division durch eine von null verschiedene Zahl
abgeschlossen.

\begin{KeyPointBox}{Vollständigkeit der reellen Zahlen}{fundamentals_real_number_completeness}
  Jeder Punkt auf der Zahlengeraden entspricht genau einer reellen Zahl und
  umgekehrt. Anschaulich besitzt die Zahlengerade daher keine Lücken: Die
  reellen Zahlen $\R$ sind \emph{vollständig}.
  \begin{center}
    \CachedTikzFigure{content/resources/chapters/fundamentals/tikz/real_number_line.tex}
  \end{center}
   Damit bilden sie eine zentrale Zahlenmenge in den Naturwissenschaften und der Technik.
\end{KeyPointBox}

Bezüglich des Wurzelziehens sind allerdings auch die reellen Zahlen nicht abgeschlossen:
Die reelle Gleichung
\[
  x^2=-1
\]
besitzt in $\R$ keine Lösung, denn das Quadrat jeder reellen Zahl ist
nichtnegativ.



\subsubsection{Ausblick: Komplexe Zahlen}
\label{subsubsec:fundamentals_complex_numbers}

Um dennoch die Wurzel aus $-1$ ziehen zu können, führt man die imaginäre
Einheit $i$ mit der Eigenschaft
\[
  i^2=-1
\]
ein. Da $i\notin\R$ gilt, ist $i$ keine reelle Zahl. Auf diese Weise erweitert
man die reellen Zahlen $\R$ zu den sogenannten \emph{komplexen Zahlen} $\C$.

\begin{center}
  \begingroup
    \def\NumberSetStage{5}
    \CachedTikzFigure{content/resources/chapters/fundamentals/tikz/number_set_inclusions.tex}
  \endgroup
\end{center}

Eine ausführliche Einführung in die komplexen Zahlen $\C$ folgt in
Kapitel~\ref{chap:complex_numbers}.



\subsection{Intervalle}
\label{subsec:fundamentals_intervals}

\begin{DefinitionBox}{Intervall}{fundamentals_interval}
  Ein \emph{Intervall} \EnglishTerm{interval} ist eine zusammenhängende
  Teilmenge $I\subseteq\R$, die zwischen zwei Randpunkten $a<b$ liegt und alle
  reellen Zahlen dazwischen enthält.

  \begin{center}
    \CachedTikzFigure{content/resources/chapters/fundamentals/tikz/interval_definition_number_line.tex}
  \end{center}
\end{DefinitionBox}

Reelle Randpunkte können zum Intervall gehören oder ausgeschlossen sein:\par

\renewcommand{\arraystretch}{1.35}
\begin{tabularx}{\linewidth}{@{}l >{\centering\arraybackslash}X l@{}}
  \toprule
  \textbf{Intervalltyp} & \textbf{Mengenschreibweise} & \textbf{Randpunkte} \\
  \midrule
  abgeschlossen
    & $[a,b]=\{x\in\R\mid a\,\Cloze{0.5cm}{\leq}\,x\,\Cloze{0.5cm}{\leq}\,b\}$
    & $a,b$ enthalten \\
  offen
    & $(a,b)=\{x\in\R\mid a\,\Cloze{0.5cm}{<}\,x\,\Cloze{0.5cm}{<}\,b\}$
    & $a,b$ nicht enthalten \\
  links abgeschlossen
    & $[a,b)=\{x\in\R\mid a\,\Cloze{0.5cm}{\leq}\,x\,\Cloze{0.5cm}{<}\,b\}$
    & nur $a$ enthalten \\
  rechts abgeschlossen
    & $(a,b]=\{x\in\R\mid a\,\Cloze{0.5cm}{<}\,x\,\Cloze{0.5cm}{\leq}\,b\}$
    & nur $b$ enthalten \\
  \bottomrule
\end{tabularx}
\renewcommand{\arraystretch}{1}

\begin{RemarkBox}[remark title=Alternative Intervallschreibweise]
  Für offene und halboffene Intervalle ist im deutschsprachigen Raum auch
  folgende Schreibweise üblich:
  \[
    (a,b)=]a,b[,
    \qquad
    [a,b)=[a,b[,
    \qquad
    (a,b]=]a,b].
  \]
\end{RemarkBox}


\begin{KeyPointBox}{Uneigentliche Randpunkte}{fundamentals_infinite_interval_endpoints}
  Liegt ein Randpunkt im Unendlichen, spricht man von einem uneigentlichen Randpunkt.
  Da $\pm \infty$ keine reellen Zahlen sind, gehören sie daher nie zum Intervall.
  An einem uneigentlichen Randpunkt (im Unendlichen) ist ein Intervall also stets offen.

  \begin{center}
    \renewcommand{\arraystretch}{1.45}
    \begin{tabularx}{0.86\linewidth}{@{}>{\centering\arraybackslash}X >{\centering\arraybackslash}X@{}}
      $(-\infty,b]$\quad\textcolor{green!50!black}{\textbf{richtig}}
      & $[a,\infty)$\quad\textcolor{green!50!black}{\textbf{richtig}} \\
      $[-\infty,b]$\quad\textcolor{red!75!black}{\textbf{falsch}}
      & $(a,\infty]$\quad\textcolor{red!75!black}{\textbf{falsch}}
    \end{tabularx}
    \renewcommand{\arraystretch}{1}
  \end{center}
\end{KeyPointBox}

